\scanpage{129}
% archive foliation: 128
Supposons $X$ connexe et distinguons les 3 cas, suivant le type d'isomorphie du rev. universel de $X$.

$1^\circ)$ Cas $S^2=\mathbb P^1_{\mathbb C}(\mathbb C)$. Alors $X\unsure{\simeq}\mathbb P^1_{\mathbb C}$ est lui-même une droite projective (car, sur une droite projective, tout autom. analytique --- et \unsure{on donne aucune condition} --- a un pt fixe au moins\ldots), $\widetilde X=X\times X\inserted{=P_X}$ avec connexion triviale, $p$ est l'identité.

$2^\circ)$ Cas $\mathbb E^1(\mathbb C)=\mathbb C$. Alors $X$ est soit \deleted{isomorphe} une courbe elliptique (cas $X$ comp. rég.) soit isomorphe à $\mathbb C$, ou à $\mathbb C^*$ (cela provient du fait que les \deleted{sous-groupes discrets}\inserted{autom. sans pt. fixe} de \deleted{$\unsure{\mathbb R}\text{ les}$} $\mathbb E^1_{\mathbb C}$, \inserted{analytiques comme de $\mathbb E^1_{\mathbb C}$, c'est kifkif}) \deleted{de $\operatorname{Aff}(1,\mathbb C)\simeq\mathbb C^*\cdot\mathbb C$ (\unsure{pour des raisons de comp.})} \inserted{sont des translations}, et les sous-groupes discrets \deleted{\unsure{contenus dans le sous-groupe des translations}} du groupe des translations $\mathbb C$ sont de rang $0$ (\unsure{point|quotient} $\mathbb C$), ou $1$ (\unsure{point|quotient} $\mathbb C^*$) ou $2$ (\unsure{point|quotient} courbe elliptique). On voit que $X$ est encore can. muni d'une structure algébrique (tout autom. analytique d'une courbe elliptique, ou d'un $\mathbb E^1_{\mathbb C}$, ou de $\mathbb G_{m\mathbb C}$, étant algébrique).

\margin{La cat. (pour iso) des courbes analytiques du type $2^\circ$ est équivalente (par le foncteur $X\mapsto X^{\mathrm{an}}$) à la catégorie des torseurs sous \deleted{\unsure{groupes linéaires}} des groupes algébriques sur $\mathbb C$ connexes de dim 1 (\unsure{comm. $\Rightarrow$ donc}, $\mathbb G_a$, $\mathbb G_m$, ou une VA de dim 1), et \unsure{can. on peut munir une telle courbe d'une structure de torseur sur $\mathbb C$}, un $\mathbb G_a$-torseur, un $\mathbb G_m$-torseur, ou une VA de dim 1).}

Ici $\widetilde X=P_X-\sigma'$, où $\sigma'$ est une \uline{section horizontale} de $P_X$ ; donc, \deleted{$P_X$ est}\inserted{$P_X$ étant} muni des deux sections analytiques $\sigma,\sigma'$ disjointes, on peut considérer $P_X$ comme clôture projective d'un fibré vectoriel analytique $V$, qui n'est autre que $\widetilde X$ comme esp. analytique sur $X$ --- mais \deleted{\unsure{muni au fond d'un fibré}}\inserted{pourvu d'une structure vectorielle}. On observera que la connexion de $V$ \deleted{\unsure{ne lui}} \unsure{n'honore} pas la structure vectorielle. Mais on peut
